% Irodov Part IV §4.2 independent original 110--125.
% User PDF pp.179--181, printed pp.170--172, answers pp.318--319.
\electricoscproblem{110}{170}
某阻尼电振荡回路品质因数为 $Q=50$、无阻尼固有振荡频率
$\nu_0=5.5\,\mathrm{kHz}$。
问储存的振荡能量下降 $\eta=2.0$ 倍需要多久？
\exsolution
弱阻尼时电流、电压振幅随 $e^{-\beta t}$ 衰减，能量包络随 $e^{-2\beta t}$ 衰减。
品质因数 $Q\simeq\omega_0/(2\beta)$，
$\omega_0=2\pi\nu_0$，
所以能量下降倍数
$\eta=e^{2\beta t}=e^{\omega_0t/Q}$。
故
\[
\boxed{t\simeq\frac{Q\ln\eta}{2\pi\nu_0}
=\frac{50\ln2}{2\pi(5500)}
\simeq1.0\,\mathrm{ms}}.
\]
\excheck 能量衰减指数是\emph{振幅}衰减指数的两倍；若忽略该因子，结果会多一倍。

\electricoscproblem{111}{170}
电振荡电路含一只漏电的电容器，电容为 $C$，其漏电电阻为 $R$；与之连接的线圈自感为 $L$，线圈及导线电阻可忽略。
求：(a) 阻尼角频率；(b) 品质因数。
\exsolution
设电容两端电压为 $V$。电感支路电流为 $I_L$，
漏电流 $I_R=V/R$，
节点 KCL
$I_L+V/R+C\dot V=0$，
电感满足 $V=L\dot I_L$（取一致有向端子）。
对 KCL 求导并消去 $I_L$：
\[
\ddot V+\frac1{RC}\dot V+\frac1{LC}V=0.
\]
因此
\[
\boxed{\beta=\frac1{2RC}},\qquad
\boxed{\omega_d=\sqrt{\frac1{LC}-\frac1{4R^2C^2}}}.
\]
按 $Q=\omega_d/(2\beta)$，
\[
\boxed{Q=RC\omega_d
=\frac12\sqrt{\frac{4R^2C}{L}-1}}.
\]
\excheck 要有阻尼振荡必须 $4R^2C>L$；这个 $R$ 是\emph{跨电容漏电支路}，不同于串联 RLC 的线圈电阻。

\electricoscproblem{112}{170}
振荡回路 $C=2.0\,\mathrm{\mu F}$、
$L=5.0\,\mathrm{mH}$，
为维持电容电压幅度 $V_m=1.0$V 的无衰减振荡，平均需输入功率
$P=0.10\,\mathrm{mW}$。阻尼很小，求品质因数。
\exsolution
回路最大存储能量
$W=\tfrac12CV_m^2$，
每周损失能量约 $P(2\pi/\omega_0)$，
$Q=2\pi W/(\text{每周耗能})=\omega_0W/P$。
因 $\omega_0=(LC)^{-1/2}$，
\[
\boxed{Q=\frac{V_m^2}{2P}\sqrt{\frac CL}}
=\frac{1}{2(10^{-4})}\sqrt{\frac{2\times10^{-6}}{5\times10^{-3}}}
=\boxed{100}.
\]
\excheck $P$ 必须以瓦特而不是毫瓦特代入；
$Q$ 为无量纲。

\electricoscproblem{113}{170}
电振荡回路中的有功电阻 $R=0.45\,\Omega$，
要维持电流幅度
$I_m=30\,\mathrm{mA}$ 的无衰减正弦振荡，
需要输入多大平均功率？
\exsolution
电流 $I=I_m\cos\omega t$，
焦耳热平均功率
$\overline P=R\langle I^2\rangle$
且 $\langle\cos^2\rangle=1/2$，
\[
\boxed{\overline P=\frac12RI_m^2}
=\frac12(0.45)(0.030)^2\,\mathrm W
=\boxed{2.0\times10^{-4}\,\mathrm W
=0.20\,\mathrm{mW}}.
\]
\excheck 书后答案将相同解析公式后标成“20 mW”，与 $30$mA、$0.45\Omega$ 直接代入差100倍。保持原版数据并登记勘误。

\electricoscproblem{114}{170}
振荡回路 $C=1.2\,\mathrm{nF}$、
$L=6.0\,\mathrm{\mu H}$、$R=0.50\,\Omega$，
要维持电容电压幅度
$V_m=10$V 的无衰减振荡，所需平均输入功率为多少？
\exsolution
弱阻尼时电感电流幅度
$I_m=\omega_0CV_m=
V_m\sqrt{C/L}$，
损耗平均值
$\overline P=\tfrac12 R I_m^2$。
因此
\[
\boxed{\overline P=\frac{RCV_m^2}{2L}}
=\frac{(0.50)(1.2\times10^{-9})(10)^2}
{2(6.0\times10^{-6})}
=\boxed{5.0\,\mathrm{mW}}.
\]
\excheck 这里 $V_m$ 是振幅；不必先换算有效值，只需显式保留因子1/2。

\electricoscproblem{115}{170}
图4.30：电容 $C$、线圈自感 $L$ 以及电阻 $R$ 分别连接在同一上、下两端点之间（并联）。
求自由阻尼振荡角频率，以及发生振荡对 $C,L,R$ 的约束。
\exsolution
图中为\emph{并联}的 L、C、R 三支路。
令两端电压为 $V$，
节点电流守恒
$C\dot V+V/R+I_L=0$。
电感关系 $L\dot I_L=V$；
对 KCL 求导：
\[
\ddot V+\frac1{RC}\dot V+\frac1{LC}V=0.
\]
因此
\[
\boxed{\omega_d=
\sqrt{\frac1{LC}-\frac1{4R^2C^2}}},
\]
仅当
\[
\boxed{R>\frac12\sqrt{\frac LC}}
\]
时有欠阻尼振荡。
\excheck 书后对并联电阻的门限是\emph{大电阻}更容易振荡；不能错套串联电阻 $R<2\sqrt{L/C}$ 的门限。

\electricoscproblem{116}{170}
图4.31有两种电振荡线路，电容均为 $C$。
(a) 两只各具自感 $L_1,L_2$、电阻 $R_1,R_2$ 的线圈与电容并联，线圈之间互感忽略；
(b) 一只自感 $L$、电阻 $R$ 的线圈与电容并联。
求使两种自由振荡的频率与阻尼相同所需的参数关系。
\exsolution
对(a)中总电感支路，在\emph{低损耗或两支路阻尼比率相容}时，可用并联电感与阻尼等效模型表示。
两自感的并联磁能等效给
\[
\boxed{\frac1L=\frac1{L_1}+\frac1{L_2}}.
\]
书后采用并联损耗的等效电导加法：
\[
\boxed{\frac1R=\frac1{R_1}+\frac1{R_2}}.
\]
在该等效模型下，
$\omega_0=1/\sqrt{LC}$、
$\beta=1/(2RC)$ 相同，因而阻尼自由频率相同。
\excheck 对图(a)中两条真正的\emph{串联 RL 支路}，
其总导纳是
$(R_1+sL_1)^{-1}+(R_2+sL_2)^{-1}$，
并不对任意参数严格等于
$(R+sL)^{-1}$。
若要求完整暂态严格等效，还需对各支路的 $R_i/L_i$ 施加匹配条件；本题书末的两条倒数关系只能在其简化电路模型下直接使用。

\electricoscproblem{117}{170}
串联 RLC 回路电容原先充电至 $V_0$，
合闸瞬间取 $t=0$。
回路电阻正好为临界阻尼值。
求随后电流 $I(t)$ 和峰值 $I_{\max}$。
\exsolution
临界阻尼
$R_{\rm cr}=2\sqrt{L/C}$，
$\beta=R/(2L)=1/\sqrt{LC}$。
电荷满足
$\ddot q+2\beta\dot q+\beta^2q=0$，
初始 $q(0)=CV_0$、$\dot q(0)=0$，
故
$q(t)=CV_0(1+\beta t)e^{-\beta t}$。
放电电流
$I=-\dot q=CV_0\beta^2 t\,e^{-\beta t}$，即
\[
\boxed{I(t)=\frac{V_0}{L}\,t\,
e^{-t/\sqrt{LC}}}.
\]
极大点 $t_m=1/\beta=\sqrt{LC}$，
\[
\boxed{I_{\max}=\frac{V_0}{e}\sqrt{\frac CL}}.
\]
\excheck 初始电流为零，$I$ 先增后减，始终不发生往复振荡，符合临界阻尼性质。

\electricoscproblem{118}{171}
电感 $L$、有功电阻 $R$ 的线圈在 $t=0$ 接入
$V(t)=V_m\cos\omega t$ 的正弦电压源。
起初无电流，求开通后 $I(t)$。
\exsolution
方程
$L\dot I+RI=V_m\cos\omega t$。
稳态项
$I_s=I_m\cos(\omega t-\varphi)$，
其中
$I_m=V_m/\sqrt{R^2+\omega^2L^2}$、
$\tan\varphi=\omega L/R$。
齐次项 $I_h=Ae^{-Rt/L}$；
由 $I(0)=0$ 有
$A=-I_m\cos\varphi$。
所以
\[
\boxed{I(t)=\frac{V_m}{\sqrt{R^2+\omega^2L^2}}
\left[\cos(\omega t-\varphi)
-\cos\varphi\,e^{-Rt/L}\right]},
\quad
\varphi=\arctan\frac{\omega L}{R}.
\]
\excheck $I(0)=0$；电感导致相位\emph{滞后}而非超前。

\electricoscproblem{119}{171}
电容 $C$ 与电阻 $R$ 串联，在 $t=0$ 接入
$V(t)=V_m\cos\omega t$ 的电源。
求开通后回路电流 $I(t)$，
并区分不同初始电容电压条件。
\exsolution
电路方程
$RI+q/C=V_m\cos\omega t$，
$I=\dot q$。
稳态电流
\[
I_s=I_m\cos(\omega t+\theta),\quad
I_m=\frac{V_m}{\sqrt{R^2+1/(\omega C)^2}},
\quad
\tan\theta=\frac1{\omega RC}.
\]
总电流
$I(t)=I_s+A e^{-t/(RC)}$，
由开通瞬间的电容电压 $V_C(0)=V_{C0}$，
$I(0)=(V_m-V_{C0})/R$，故
\[
\boxed{I(t)=I_m\cos(\omega t+\theta)
+\left[\frac{V_m-V_{C0}}R-I_m\cos\theta\right]
e^{-t/(RC)}}.
\]
若电容初始不带电 $V_{C0}=0$，直接代入即可。
\textbf{书末约定：}书末给出的形式
$I=I_m[\cos(\omega t-\varphi)-
\cos\varphi e^{-t/(RC)}]$（$\varphi=-\theta$）
对应 $I(0)=0$，即预先令
$V_{C0}=V_m$，而\emph{不是}初始未充电状态。
\excheck 原题未指定初始电容电压，故单一暂态解不唯一；本稿先给包含 $V_{C0}$ 的完整解，再说明书末隐含初值。

\electricoscproblem{120}{171}
长单层密绕螺线管绕线电阻率为 $\rho$，
单位轴向长度匝数 $n$，
螺线管截面半径为 $a$，
绝缘层厚度可忽略。
施加频率 $\nu$ 的交变电压时，
求电流与电压的相位差。
\exsolution
导线直径约 $d=1/n$，
截面积 $A_w=\pi/(4n^2)$。
每单位螺线管轴向长度的导线长度
$s'=2\pi an$，所以
$R'=\rho s'/A_w=8\rho an^3$。
长螺线管单位轴长自感
$L'=\mu_0n^2\pi a^2$。
交变角频率 $\Omega=2\pi\nu$，
相位滞后
$\tan\varphi=\Omega L'/R'$，
故
\[
\boxed{\tan\varphi=
\frac{\mu_0\pi^2\nu a}{4\rho n}}.
\]
\excheck 该表达式适用长、单层、紧密绕制和忽略趋肤效应；$n$ 是每米匝数，不能当总匝数。

\electricoscproblem{121}{171}
电容 $C$ 与电阻 $R=110\,\Omega$ 串联，
接到电压幅度 $V_m=110$V 的交流源，
测得稳态电流幅度 $I_m=0.50$A。
求电流与源电压的相位差。
\exsolution
阻抗模
$|Z|=V_m/I_m=220\,\Omega$。
串联 RC 中
$\cos\varphi=R/|Z|=110/220=1/2$。
因此电流相对电源电压\emph{超前}
\[
\boxed{\varphi=\arccos(1/2)=60^\circ}.
\]
\excheck 电路为容性，电流必须超前；若以电压相对电流的相位定义，则记为 $-60^\circ$。

\electricoscproblem{122}{171}
图4.32给出最简脉动滤波器：输入
$V_{\rm in}(t)=V_0(1+\cos\omega t)$，
经串联电阻 $R$ 后输出跨并联电容 $C$ 的电压。
求：(a) 稳态输出；(b) $\omega=314\,\mathrm{s^{-1}}$ 时，
输出交流幅度为输出直流分量的 $1/\eta$，
$\eta=7.0$，所需 $RC$。
\exsolution
直流稳态下电容开路，输出常值部分为 $V_0$。
交流部分采用传递函数
$H(i\omega)=1/(1+i\omega RC)$，
因此
\[
\boxed{V_{\rm out}(t)=V_0+
\frac{V_0}{\sqrt{1+(\omega RC)^2}}
\cos\left(\omega t-\arctan\omega RC\right)}.
\]
要求
$\sqrt{1+(\omega RC)^2}=\eta$，
所以
\[
\boxed{RC=\frac{\sqrt{\eta^2-1}}{\omega}
=\frac{\sqrt{48}}{314}
\simeq22\,\mathrm{ms}}.
\]
\excheck 输出直流增益为1，交流增益小于1；频率增大则滤波效果增强。

\electricoscproblem{123}{171}
图4.33(a)(b)给出两种正弦交流串联电路，
要求按外加电压与电流的相位关系作近似电压相量图。
(a) 纯电阻 $R_0$ 与具有自感 $L$、电阻 $R$ 的线圈串联；
(b) 电容 $C$ 与具有自感 $L$、电阻 $R$ 的线圈串联，
且题设 $\omega<\omega_0=1/\sqrt{LC}$。
\exsolution
取电流相量 $\bm I$ 沿实轴。
(a) 电阻 $R_0$ 上电压
$\underline V_{R_0}=IR_0$ 同相，
线圈电压
$\underline V_{\rm coil}=I(R+i\omega L)$，
总电压
\[
\boxed{\underline V=I[(R_0+R)+i\omega L]},
\quad
\tan\varphi=\frac{\omega L}{R_0+R}>0.
\]
相量图中两段水平电阻电压相加，感抗电压竖直向上，总电压在第一象限。
(b) 电容电压
$\underline V_C=-iI/(\omega C)$，
线圈电压如上；
$\omega<\omega_0$ 时
$\omega L-1/(\omega C)<0$，故
\[
\boxed{\underline V=I\left[
R+i\left(\omega L-\frac1{\omega C}\right)\right]},
\]
总电压相对电流落在第四象限；相量图中的电容电压向下长于电感向上的电压。
\excheck 用复数相量的实轴和正负虚轴描述给出了原图4.33的两个电压矢量图的全部相对方向与长度约束；不用电压有效值和峰值混算。

\electricoscproblem{124}{171}
串联 RLC 回路含
$C=22\,\mathrm{\mu F}$、线圈
$R=20\,\Omega$、$L=0.35$H，
外加电压幅度 $V_m=180$V，
角频率 $\omega=314\,\mathrm{s^{-1}}$。
求：(a) 电流幅度；(b) 电流与外电压相位差；
(c) 电容及线圈\emph{整个线圈端子}的电压幅度。
\exsolution
感抗 $X_L=\omega L\simeq109.9\,\Omega$，
容抗 $X_C=1/(\omega C)\simeq144.8\,\Omega$。
阻抗模
$Z=\sqrt{R^2+(X_L-X_C)^2}\simeq40.2\,\Omega$。
故
\[
\boxed{I_m=\frac{V_m}{Z}\simeq4.5\,\mathrm A}.
\]
以源电压相对电流的相位为 $\varphi$：
\[
\boxed{\tan\varphi=\frac{X_L-X_C}{R}},
\quad\boxed{\varphi\simeq-60^\circ},
\]
即电流超前电压约 $60^\circ$。
电容两端幅度
$\boxed{V_{Cm}=I_mX_C\simeq0.65\,\mathrm{kV}}$；
整个线圈含内阻，其端电压幅度
\[
\boxed{V_{{\rm coil},m}
=I_m\sqrt{R^2+X_L^2}\simeq0.50\,\mathrm{kV}}.
\]
\excheck 交流相量叠加不等于电压幅值的代数相加，
因此单个器件的电压幅度可大于电源幅度。

\electricoscproblem{125}{171--172}
正弦电压幅度保持不变、频率可调的电源，
向电容 $C$、纯电阻 $R$、无电阻电感 $L$ 的串联回路供电。
求令电压\emph{幅度}最大的驱动频率：
(a) 跨电容；
(b) 跨电感。
\exsolution
记 $\omega_0^2=1/(LC)$、
$\beta=R/(2L)$，驱动频率为 $\Omega$。
电流幅度
$I_m=V_m/\sqrt{R^2+(\Omega L-1/(\Omega C))^2}$。
电容电压
$V_C=I_m/(\Omega C)$，
整理平方：
\[
V_C^2=\frac{V_m^2\omega_0^4}
{(\omega_0^2-\Omega^2)^2+4\beta^2\Omega^2}.
\]
分母对 $\Omega^2$ 求极小，
\[
\boxed{\Omega_{C,\max}=
\sqrt{\omega_0^2-2\beta^2}}.
\]
电感电压
$V_L=I_m\Omega L$，对应
\[
V_L^2=\frac{V_m^2\Omega^4}
{(\omega_0^2-\Omega^2)^2+4\beta^2\Omega^2}.
\]
令 $y=\Omega^2$，对上式求导得极值
\[
\boxed{\Omega_{L,\max}
=\frac{\omega_0^2}
{\sqrt{\omega_0^2-2\beta^2}}}.
\]
\excheck 两式要求 $\omega_0^2>2\beta^2$；
高阻尼时不存在同样的内部有限频率幅度峰。
